\documentclass[11pt,letterpaper]{article}
\usepackage[utf8]{inputenc}
\usepackage{geometry}
\usepackage{setspace}
\usepackage{enumitem}

\geometry{left=1.0in,right=1.0in,top=1.0in,bottom=1.0in}
\doublespacing

\setlength{\parindent}{0pt}
\setlength{\parskip}{12pt plus 2pt minus 2pt}

\begin{document}

% =========================================================================
% REPLACEMENT ABSTRACT (37 CFR § 1.72(b) - COMMENCING ON A SEPARATE SHEET)
% =========================================================================
\thispagestyle{empty}

\begin{center}
    \textbf{\large REPLACEMENT ABSTRACT} \\
    \vspace{0.3cm}
    \textbf{Application No.: 19/639,222} \\
    \textbf{Art Unit: Pre-Exam / TC 2100} \\
    \textbf{Examiner: Andrea Bennett}
\end{center}

\vspace{0.8cm}

\begin{center}
    \textbf{A MEMORY SYSTEM FOR THE STORAGE AND RETRIEVAL OF PARAMETRIC AND SYMBOLIC DATA REPRESENTATIONS}
\end{center}

\vspace{0.5cm}

\begin{center}
    \textbf{ABSTRACT}
\end{center}

A memory system and architecture designed for the storage and retrieval of data as parametric or symbolic representations rather than as fixed-width, evaluated numerical values. The system comprises a memory controller and a storage medium (such as DRAM, flash, or ROM) where data is stored as variable-length structures representing mathematical 'recipes' or expressions. The memory controller manages the dynamic allocation of physical memory to accommodate these structures and interfaces with a parametric data bus. This architecture ensures that the precision of data is not lost when it is moved from a parametric CPU to memory, enabling the persistence of perfectly represented fractional, irrational, and large-magnitude numbers throughout a complete computing system.

\newpage

% =========================================================================
% REPLACEMENT CLAIMS (37 CFR § 1.75(h) - COMMENCING ON A SEPARATE SHEET)
% =========================================================================
\setcounter{page}{1}
\pagestyle{plain}

\begin{center}
    \textbf{\large REPLACEMENT CLAIMS} \\
    \vspace{0.3cm}
    \textbf{Application No.: 19/639,222} \\
    \textbf{Art Unit: Pre-Exam / TC 2100} \\
    \textbf{Examiner: Andrea Bennett}
\end{center}

\vspace{0.8cm}

\textbf{What is claimed is:}

1. A memory system for storing data, comprising:
\begin{itemize}[leftmargin=2.5em]
    \item[] a physical storage medium; and
    \item[] a parametric memory controller configured to:
    \begin{itemize}[leftmargin=2.5em]
        \item[--] receive a data structure representing a mathematical expression of a number;
        \item[--] dynamically allocate a variable-length block of the physical storage medium to store said data structure; and
        \item[--] retrieve the complete, unmodified data structure from the storage medium upon a read request.
    \end{itemize}
\end{itemize}

2. The memory system of claim 1, wherein the data structure is received from a parametric central processing unit over a parametric data bus.

3. The memory system of claim 1, wherein the mathematical expression represents a fractional number, and said representation is stored without loss of precision.

4. The memory system of claim 1, wherein the mathematical expression is a symbolic representation of a mathematical or physical constant.

5. The memory system of claim 1, wherein the parametric memory controller manages an addressing index to locate the variable-length blocks within the physical storage medium.

6. The memory system of claim 1, wherein the physical storage medium is volatile random-access memory (RAM).

7. The memory system of claim 1, wherein the physical storage medium is non-volatile memory.

8. The memory system of claim 7, wherein the physical storage medium is read-only memory (ROM) pre-configured with a library of data structures representing mathematical constants.

9. A method for storing data, comprising:
\begin{itemize}[leftmargin=2.5em]
    \item[] receiving, at a memory controller, a data structure representing a mathematical expression;
    \item[] determining the size of the received data structure;
    \item[] allocating a block of memory in a storage medium, wherein the size of the block is based on the determined size of the data structure;
    \item[] storing the entire data structure in the allocated block; and
    \item[] creating an index entry that maps a memory address to the location of the allocated block.
\end{itemize}

\vspace{1.0cm}
\begin{singlespace}
\noindent Respectfully submitted, \par\vspace{1.2cm}
\underline{\textbf{/Reinaldo Javier Aguiar Marcano/}} \\
\textbf{Reinaldo Javier Aguiar Marcano} \\
Pro Se Applicant \\
Date: October 9, 2026
\end{singlespace}

\end{document}
